% Part: normal-modal-logic
% Chapter: axioms-systems
% Section: normal-logics

\documentclass[../../../include/open-logic-section]{subfiles}

\begin{document}

\olfileid{nml}{prf}{nor}

\olsection{सामान्य मोडल तर्कशास्त्रे}

मोडल !!{formula}s चा प्रत्येक संच स्वयंसिद्धकांच्या
एखाद्या संचापासून सिद्ध करता येणाऱ्या
!!{formula}s चा संच म्हणून सहज वर्णन करता येत
नाही. मोडल तर्कशास्त्रे सुस्थित असावीत अशी
आपली अपेक्षा आहे. सर्वप्रथम, अभिजात विधानीय
तर्कशास्त्रात आपण !!{derive} करू शकतो ते सर्व
अजूनही !!{derivable} असावे; आता !!{formula}s
मध्ये \iftag{prvBox}{$\Box$\iftag{prvDiamond}{
    आणि~}{}}{}\iftag{prvDiamond}{$\Diamond$}{} देखील
येऊ शकतात, हे अर्थातच लक्षात घेतले पाहिजे.
यासाठी मोडल तर्कशास्त्रात सर्व सर्वतः-सत्य दृष्टांत
असावेत आणि ते विध्यात्मक अनुमानाखाली बंद असावे,
अशी अट घालतो.

\begin{defn}
  \emph{मोडल तर्कशास्त्र} म्हणजे मोडल !!{formula}s
  चा असा संच~$\Sigma$ की, तो
  \begin{enumerate}
  \item सर्व सर्वतः-सत्य सूत्रे सामावतो; आणि
  \item आदेशनाखाली बंद असतो; म्हणजे $!A \in \Sigma$
    आणि $!D_1$, \dots, $!D_n$ ही !!{formula}s
    असतील, तर
    \[
    \SSubst{!A}{\subst{!D_1}{p_1}, \dots, \subst{!D_n}{p_n}} \in \Sigma,
    \]
    \item \emph{विध्यात्मक अनुमानाखाली} बंद असतो;
    म्हणजे $!A$ आणि $!A
      \lif !B \in \Sigma$ असतील, तर $!B \in \Sigma$.
  \end{enumerate}
\end{defn}

मोडल तर्कशास्त्रांसाठी संबंधी चिन्हार्थमीमांसा
वापरायची असेल, तर सर्व मोडल प्रतिमानांत वैध
असलेली सर्व !!{formula}s त्यांत समाविष्ट असावीत,
अशीही अट हवी. \Ax{K}\iftag{prvDiamond}{ आणि~\Dual{}}{}
चे सर्व दृष्टांत !!{derivable} असतील, आणि
!!a{formula}~$!A$ !!{derivable} असेल तेव्हा
$\Box !A$ देखील तसेच असेल, तर ही अट पूर्ण होते.
या अटी पाळणाऱ्या मोडल तर्कशास्त्राला
\emph{सामान्य} म्हणतात. (असामान्य मोडल तर्कशास्त्रेही
आहेत; पण नेहमीची संबंधी प्रतिमाने त्यांच्यासाठी
पुरेशी नाहीत.)

\begin{defn}
  मोडल तर्कशास्त्र $\Sigma$ पुढील सूत्रे सामावत असेल,
  \begin{align*}
    \tag{\Ax{K}} & \Box(p \lif q) \lif (\Box p \lif \Box q),
    \iftag{prvDiamond}{\\
      \tag{\Dual} & \Diamond p \liff \lnot\Box\lnot p}{}
  \end{align*}
  आणि \emph{अनिवार्यीकरणाखाली} बंद असेल, म्हणजे
  $!A \in
  \Sigma$ असेल, तर $\Box !A \in \Sigma$,
  तर त्याला \emph{सामान्य} म्हणतात.
\end{defn}

सर्वतः-सत्य अभिव्यंजन \emph{जगातील सत्यता}
टिकवण्यासाठी पुरेसे सूक्ष्म असते; पण \Nec{} हा
नियम फक्त \emph{प्रतिमानातील सत्यता} टिकवतो
(आणि म्हणून चौकटीतील किंवा चौकटींच्या वर्गातील
वैधताही टिकवतो), हे लक्षात घ्या.

\begin{prop}\ollabel{prop:rk}
  प्रत्येक सामान्य मोडल तर्कशास्त्र \RK
  या नियमाखाली बंद असते:
  \begin{prooftree}
    \AxiomC{$!A_1 \lif (!A_2 \lif \cdots (!A_{n-1} \lif !A_n)\cdots)$}
    \RightLabel{\RK}
    \UnaryInfC{$\Box!A_1 \lif (\Box!A_2 \lif \cdots (\Box!A_{n-1}
      \lif \Box!A_n)\cdots).$}
  \end{prooftree}
\end{prop}

\begin{proof}
  $n$ वर विगमनाने सिद्ध करू. $n = 1$ असेल,
  तर हा नियम नेमका \Nec असतो, आणि प्रत्येक
  सामान्य मोडल तर्कशास्त्र \Nec खाली बंद असते.

  आता निष्कर्ष $n-1$ साठी खरा आहे असे समजा;
  तो~$n$ साठी सिद्ध करू.

  पुढील गृहीत धरा:
  \begin{align*}
  & !A_1 \lif (!A_2 \lif \cdots (!A_{n-1} \lif !A_n)\cdots) \in \Sigma
  \intertext{विगमन गृहीतकावरून मिळते:}
  & \Box!A_1 \lif (\Box!A_2 \lif \cdots \Box(!A_{n-1} \lif !A_n)\cdots)
  \in \Sigma
  \intertext{$\Sigma$ सामान्य मोडल तर्कशास्त्र असल्यामुळे
    त्यात~$\Ax{K}$ चे सर्व दृष्टांत आहेत; विशेषतः:}
  & \Box(!A_{n-1} \lif !A_n) \lif (\Box!A_{n-1} \lif \Box!A_n) \in \Sigma
  \intertext{विध्यात्मक अनुमान आणि योग्य सर्वतः-सत्य दृष्टांत
    वापरून मिळते:}
  & \Box!A_1 \lif (\Box!A_2 \lif \cdots (\Box!A_{n-1}
  \lif \Box!A_n)\cdots) \in \Sigma.
  \end{align*}
\end{proof}

\begin{prop}\ollabel{prop:notDiamondBot}
  प्रत्येक सामान्य मोडल तर्कशास्त्र $\Sigma$
  हे~$\lnot\Diamond\lfalse$ सामावते.
\end{prop}

\begin{prob}
  \olref[nml][prf][nor]{prop:notDiamondBot} सिद्ध करा.
\end{prob}

\begin{prop}
  $!A_1$, \dots, $!A_n$ ही !!{formula}s असू द्या.
  मग $!A_1$, \dots, $!A_n$ यांचे सर्व दृष्टांत
  सामावणारे लघुतम सामान्य मोडल तर्कशास्त्र $\Sigma$
  अस्तित्वात असते.
\end{prop}

\begin{proof}
  $!A_1$, \dots, $!A_n$ दिल्यावर
  $!A_1$, \dots, $!A_n$ यांचे सर्व दृष्टांत सामावणाऱ्या
  सर्व सामान्य मोडल तर्कशास्त्रांचा
  छेद~$\Sigma$ म्हणून ठरवा. अशा तर्कशास्त्रांचा
  वर्ग रिकामा नाही: सर्व !!{formula}s चा संच
  $\Frm[L]$ हेही असे सामान्य मोडल तर्कशास्त्र आहे.
\end{proof}

\begin{defn}
$!A_1$, \dots, $!A_n$ सामावणाऱ्या लघुतम
सामान्य मोडल तर्कशास्त्राला \emph{मोडल प्रणाली}
म्हणतात; ते $\Log{K} !A_1 \dots
!A_n$ असे दर्शवतात. लघुतम सामान्य मोडल
तर्कशास्त्र~\Log{K} असे दर्शवतात.
\end{defn}

\end{document}
